% !PRINTED 467
\weberpage{467}

Nun haben $G(t)$ und $g(t)$ eine Wurzel gemein. Da aber $g(t)$ irreducibel ist, so muss $G(t)$ durch $g(t)$ theilbar sein; also sind auch alle Wurzeln von $g(t)$ in $N$ enthalten, d.\,h. $N$ ist ein Normalkörper, w.\,z.\,b.\,w.

Jede Gleichung $g(t) = 0$ heisst eine \gesperrt{Galois'sche Resolvente} der Gleichung $F(x) = 0$, und eine \gesperrt{Galois'sche Resolvente} ist also durch folgende Bestimmung definirt:

\gesperrt{Eine Gleichung $g(t) = 0$ ist eine Galois'sche Resolvente einer gegebenen Gleichung $F(x) = 0$ in $\Omega$, wenn 1) $g(t)$ irreducibel ist, wenn 2) alle Wurzeln von $F(x)$ rational durch eine Wurzel $\varrho$ von $g(t)$ ausdrückbar sind, und 3) eine Wurzel von $g(t)$ rational durch die Wurzeln von $F(x)$ ausdrückbar ist.}

Denn nach 2) ist $N$ in $\Omega(\varrho)$ enthalten, und nach 3) ist einer der mit $\Omega(\varrho)$ conjugirten Körper, $\Omega(\varrho_1)$, in $N$ enthalten. Die Grade von $N$ und $\Omega(\varrho)$ können also nicht verschieden sein und folglich ist $N = \Omega(\varrho)$.

\gesperrt{Jede Galois'sche Resolvente ist eine Normalgleichung, und eine Normalgleichung ist ihre eigene Galois'sche Resolvente.}

\begin{center}
\S.\,146.

\gesperrt{Die Substitutionen eines Normalkörpers.}
\end{center}

Es sei jetzt $\Omega(\varrho)$ ein Normalkörper $\mu^{\mathrm{ten}}$ Grades und $\varrho$ eine seiner primitiven Zahlen, ferner
\weq{1}{g(t) = 0}
die irreducible Gleichung $\mu^{\mathrm{ten}}$ Grades, deren eine Wurzel $t = \varrho$ ist. Die zu $\varrho$ conjugirten Elemente seien
\weq{2}{\varrho,\ \varrho_1,\ \varrho_2 \ldots \varrho_{\mu-1}.}
Da nach der Definition des Normalkörpers die Grössen (2) alle in $\Omega(\varrho)$ enthalten sind, so können wir $\Theta_1(t), \Theta_2(t) \ldots \Theta_{\mu-1}(t)$ als ganze rationale Functionen in $\Omega$, höchstens vom Grade $\mu - 1$, so bestimmen, dass
\weq{3}{\varrho_1 = \Theta_1(\varrho),\ \varrho_2 = \Theta_2(\varrho) \ldots \varrho_{\mu-1} = \Theta_{\mu-1}(\varrho)}
%CONT>>
